\subsection{Groups Acting on Themselves by Conjugation---The Class Equation}

Let $G$ be a group.

\begin{specialexercise}
	Suppose $G$ has a left action on a set $A$, denoted by $g \cdot a$ for all $g \in G$ and $a \in A$. Denote the corresponding right action on $A$ by $a \cdot g$. Prove that the (equivalence) relations $\sim$ and $\sim'$ defined by
	\[ a \sim b \quad \text{if and only if} \quad a = g \cdot b \quad \text{ for some } g \in G \]
	and
	\[ a \sim' b \quad \text{if and only if} \quad a = b \cdot g \quad \text{ for some } g \in G \]
	are the same relation (i.e., $a \sim b$ if and only if $a \sim' b$).
\end{specialexercise}

\begin{sol}
	If $a \sim b$, there exists $g \in G$ such that $a = g \cdot b$. Let $h = g\inv$. Then $b = h \cdot a$, so $a = b \cdot g$. Thus, $a \sim' b$. Conversely, if $a \sim' b$, there exists $g \in G$ such that $a = b \cdot g$. Let $h = g\inv$. Then $b = h \cdot a$, so $a = g \cdot b$. Thus, $a \sim b$. Therefore, the relations $\sim$ and $\sim'$ are the same.
\end{sol}

\begin{exercise}
	[4.3.2] Find all conjugacy classes and their sizes in the following groups:
	\begin{subexercises}
		\item $D_8$ \quad (b)~~$Q_8$ \quad (c)~~$A_4$
	\end{subexercises}
\end{exercise}

\begin{solenum}
	\item Discussed in the text, the conjugacy classes of $D_8$ are $\{1\}$, $\{r^2\}$, $\{r, r^3\}$, $\{s, sr^2\}$, and $\{sr, sr^3\}$ with sizes 1, 1, 2, 2, and 2 respectively.
	\item Again discussed in the text, the conjugacy classes of $Q_8$ are $\{1\}$, $\{-1\}$, $\{i, -i\}$, $\{j, -j\}$, and $\{k, -k\}$ with sizes 1, 1, 2, 2, and 2 respectively.
	\item We proceed similarly as the text. The possible cycle types are $1, (1\ 2\ 3)$, and $(1\ 2)(3\ 4)$.
	
	For $(1\ 2\ 3)$, we have $C_{A_4}((1\ 2\ 3)) = \gen{(1\ 2\ 3)}$ since the only permutation that fixes 1, 2, and 3 is the identity. Therefore, the conjugacy class of $(1\ 2\ 3)$ has size $12/3 = 4$. However, there are 8 3-cycles in $A_4$, so there exists a 3-cycle $\sigma \in A_4$ not in the conjugacy class of $(1\ 2\ 3)$. By similar reasoning, the conjugacy class of $\sigma$ also has size 4. There are then 2 conjugacy classes of size 4 corresponding to the 3-cycles.

	For $(1\ 2)(3\ 4)$, it is trivial to calculate that the remaining double transpositions are in its conjugacy class. Therefore, the conjugacy class of $(1\ 2)(3\ 4)$ has size 3.
\end{solenum}

\newpage

\begin{exercise}
	Find all conjugacy classes and their sizes in the following groups:
	\begin{subexercises}
		\item $Z_2 \times S_3$ \quad (b)~~$S_3 \times S_3$ \quad (c)~~$Z_3 \times A_4$
	\end{subexercises}
\end{exercise}

\begin{sol}
	We first prove a (somewhat) trivial idea: if $G$ and $H$ are groups, let $g \in G$ and $h \in H$. Let $\kk_g$ and $\kk_h$ be the conjugacy classes of $g$ in $G$ and $h$ in $H$ respectively. Then $\kk_{(g,h)} = \kk_g \times \kk_h$ is the conjugacy class of $(g,h)$ in $G \times H$. To see this, note that for any $(x,y) \in G \times H$, we have
	\[(x,y)(g,h)(x,y)\inv = (xgx\inv, yhy\inv) \in \kk_g \times \kk_h.\]
	Conversely, for any $(g', h') \in \kk_g \times \kk_h$, there exist $x \in G$ and $y \in H$ such that $g' = xgx\inv$ and $h' = yhy\inv$. Therefore, $(g', h') = (x,y)(g,h)(x,y)\inv$, so $(g', h')$ is in the conjugacy class of $(g,h)$ in $G \times H$. This result shows two important things: the conjugacy classes of $G \times H$ are precisely the products of the conjugacy classes of $G$ and $H$, and the size of the conjugacy class of $(g,h)$ in $G \times H$ is the product of the sizes of the conjugacy classes of $g$ in $G$ and $h$ in $H$. We now use this to answer the question.
	\begin{subexercises}
		\item Let $Z_2 = \gen x$. The conjugacy classes $\kk_1$ and $\kk_x$ of $Z_2$ have sizes 1 and 1 respectively. For $S_3$, the partitions of 3 are 3 1-cycles, 2-cycle and 1-cycle, and a 3-cycle with representatives $1$, $(1\ 2)$, and $(1\ 2\ 3)$ respectively. Since $C_{S_3}((1\ 2)) = \gen{(1\ 2)}$ and $C_{S_3}((1\ 2\ 3)) = \gen{(1\ 2\ 3)}$, then the conjugacy classes $\kk_1$, $\kk_{(1\ 2)}$, and $\kk_{(1\ 2\ 3)}$ of $S_3$ have sizes 1, 3, and 2 respectively. Therefore, $Z_2 \times S_3$ has 6 conjugacy classes of size 1, 3, 2, 1, 3, and 2 respectively.
		\item By the above, $S_3$ has 3 conjugacy classes of sizes 1, 3, and 2 respectively. Then $S_3 \times S_3$ has 9 conjugacy classes with sizes $1, 3, 2, 3, 9, 6, 2, 6,$ and $4$ respectively.
		\item $Z_3$ has 3 conjugacy classes of size 1 each. From \exref{4.3.2}, we know that $A_4$ has 4 conjugacy classes of sizes 1, 4, 4, and 3. Therefore, $Z_3 \times A_4$ has 12 conjugacy classes with sizes $1, 4, 4, 3, 1, 4, 4, 3, 1, 4, 4,$ and $3$ respectively. \qh
	\end{subexercises}
\end{sol}

\begin{exercise}
	Prove that if $S \subseteq G$ and $g \in G$ then $gN_G(S)g^{-1} = N_G(gSg^{-1})$ and $gC_G(S)g^{-1} = C_G(gSg^{-1}).$
\end{exercise}

\begin{sol}
	Suppose $x \in gN_G(S)g\inv$, and consider $xgSg\inv x\inv$. Since $x = gng\inv$ for some $n \in N_G(S)$, then we have $gng\inv gSg\inv gn\inv g\inv = g n S n\inv g\inv = g S g\inv$, so that $x \in N_G(gSg\inv)$. Conversely, suppose $x \in N_G(gSg\inv)$. Then $xgSg\inv x\inv = gSg\inv$. Letting $n = g\inv x g$, we have $n S n\inv = S$, so $n \in N_G(S)$ and hence $x \in gN_G(S)g\inv$. Therefore, $gN_G(S)g\inv = N_G(gSg\inv)$.

	Suppose $x \in gC_G(S)g\inv$, and consider $xsx\inv$ for some $s \in gSg\inv$. Since $x = gng\inv$ for some $n \in C_G(S)$, then we have $gng\inv s g n\inv g\inv = g n (g\inv s g) n\inv g\inv = g (g\inv s g) g\inv = s$, so that $x \in C_G(gSg\inv)$. Conversely, suppose $x \in C_G(gSg\inv)$. Then $xsx\inv = s$ for all $s \in gSg\inv$. Letting $n = g\inv x g$, we have $n t n\inv = t$ for all $t \in S$, so $n \in C_G(S)$ and hence $x \in gC_G(S)g\inv$. Therefore, $gC_G(S)g\inv = C_G(gSg\inv)$.
\end{sol}

\begin{exercise}
	If the center of $G$ is of index $n$, prove that every conjugacy class has at most $n$ elements.
\end{exercise}

\begin{sol}
	Let $\kk_g$ be the conjugacy class of some $g \in G$. By Proposition 4.6, we have $|\kk_g| = |G : C_G(g)|$. Since $Z(G) \leq C_G(g)$ for all $g \in G$, Lagrange's Theorem shows that $|Z(G)|$ divides $|C_G(g)|$. Therefore, $|G : C_G(g)|$ divides $|G : Z(G)| = n$. Hence, $|\kk_g| \leq n$.
\end{sol}

\newpage

\begin{specialexercise}
	Assume $G$ is a non-abelian group of order $15$. Prove that $Z(G)=1$. Use the fact that $\langle g \rangle \le C_G(g)$ for all $g \in G$ to show that there is at most one possible class equation for $G$. [Use \exref{3.1.36}.]
\end{specialexercise}

\begin{sol}
	Recall that $Z(G) \leq G$. Since $G$ is non-abelian, then $Z(G) \neq G$. Assume $1 < Z(G) < 15$. By Lagrange's Theorem, $Z(G)$ has order 3 or 5. If $|Z(G)| = 3$, then $|G/Z(G)| = 5$ and $G/Z(G) \cong Z_5$. If $|Z(G)| = 5$, then $|G/Z(G)| = 3$ and $G/Z(G) \cong Z_3$. In either case, $G/Z(G)$ is cyclic by Corollary 3.10, and \exref{3.1.36} implies that $G$ is abelian, a contradiction. Therefore, $Z(G) = 1$.

	Now suppose $g \in G$ is nonidentity. Since $\langle g \rangle \leq C_G(g)$, then $|C_G(g)|$ is 3, 5, or 15 by Lagrange's Theorem. If $|C_G(g)| = 15$, then $C_G(g) = G$, so $g \in Z(G)$, a contradiction. Therefore, $|C_G(g)|$ is 3 or 5 for all nonidentity $g \in G$. By Proposition 4.6, we have $|\kk_g| = |G : C_G(g)|$, so that $|\kk_g|$ is 5 or 3 respectively. The identity element forms a conjugacy class of size 1, and the class equation must involve a summation of sizes 3 and 5 that adds to 14. The only such combination is three classes of size 3 and one class of size 5. Therefore, the class equation of $G$ is $15 = 1 + 3 + 3 + 3 + 5$.
\end{sol}

\begin{exercise}
	For $n=3,4,6,$ and $7$ make lists of the partitions of $n$ and give representatives for the corresponding conjugacy classes of $S_n$.
\end{exercise}

\begin{sol}
	$n = 3$:
	\[\begin{array}{c|c}
		\textbf{Partition of 3} & \textbf{Representative of Cycle Type} \\
		\hline
		1, 1, 1 & 1 \\
		2, 1 & (1\ 2) \\
		3 & (1\ 2\ 3)
	\end{array}\]
	$n = 4$:
	\[\begin{array}{c|c}
		\textbf{Partition of 4} & \textbf{Representative of Cycle Type} \\
		\hline
		1, 1, 1, 1 & 1 \\
		2, 1, 1 & (1\ 2) \\
		2, 2 & (1\ 2)(3\ 4) \\
		3, 1 & (1\ 2\ 3) \\
		4 & (1\ 2\ 3\ 4)
	\end{array}\]
	$n = 6$:
	\[\begin{array}{c|c}
		\textbf{Partition of 6} & \textbf{Representative of Cycle Type} \\
		\hline
		1, 1, 1, 1, 1, 1 & 1 \\
		2, 1, 1, 1, 1 & (1\ 2) \\
		2, 2, 1, 1 & (1\ 2)(3\ 4) \\
		2, 2, 2 & (1\ 2)(3\ 4)(5\ 6) \\
		3, 1, 1, 1 & (1\ 2\ 3) \\
		3, 2, 1 & (1\ 2\ 3)(4\ 5) \\
		3, 3 & (1\ 2\ 3)(4\ 5\ 6) \\
		4, 1, 1 & (1\ 2\ 3\ 4) \\
		4, 2 & (1\ 2\ 3\ 4)(5\ 6) \\
		5, 1 & (1\ 2\ 3\ 4\ 5) \\
		6 & (1\ 2\ 3\ 4\ 5\ 6)
	\end{array}\]
	$n = 7$:
	\[\begin{array}{c|c}
		\textbf{Partition of 7} & \textbf{Representative of Cycle Type} \\
		\hline
		1, 1, 1, 1, 1, 1, 1 & 1 \\
		2, 1, 1, 1, 1, 1 & (1\ 2) \\
		2, 2, 1, 1, 1 & (1\ 2)(3\ 4) \\
		2, 2, 2, 1 & (1\ 2)(3\ 4)(5\ 6) \\
		3, 1, 1, 1, 1 & (1\ 2\ 3) \\
		3, 2, 1, 1 & (1\ 2\ 3)(4\ 5) \\
		3, 2, 2 & (1\ 2\ 3)(4\ 5)(6\ 7) \\
		3, 3, 1 & (1\ 2\ 3)(4\ 5\ 6) \\
		4, 1, 1, 1 & (1\ 2\ 3\ 4) \\
		4, 2, 1 & (1\ 2\ 3\ 4)(5\ 6) \\
		4, 3 & (1\ 2\ 3\ 4)(5\ 6\ 7) \\
		5, 1, 1 & (1\ 2\ 3\ 4\ 5) \\
		5, 2 & (1\ 2\ 3\ 4\ 5)(6\ 7) \\
		6, 1 & (1\ 2\ 3\ 4\ 5\ 6) \\
		7 & (1\ 2\ 3\ 4\ 5\ 6\ 7)
	\end{array} \qh\]
\end{sol}

\begin{exercise}
	Prove that $Z(S_n)=1$ for all $n \ge 3$.
\end{exercise}

\begin{sol}
	Suppose $\sigma \in Z(S_n)$ for some $n \geq 3$. Then $\sigma$ commutes with every element of $S_n$. In particular, $\sigma$ commutes with all transpositions (the set of which generate $S_n$). Let $(a\ b)$ be any transposition in $S_n$. Then we have $\sigma(a\ b)\sigma\inv = (a\ b)$. This implies that $(\sigma(a)\ \sigma(b)) = (a\ b)$, so $\sigma(a) = a$ and $\sigma(b) = b$. Since $a$ and $b$ were arbitrary, $\sigma$ fixes every element of $\{1, 2, \ldots, n\}$. Therefore, $\sigma$ is the identity permutation. Hence, $Z(S_n) = \{1\}$ for all $n \geq 3$.
\end{sol}

\begin{exercise}
	Show that $|C_{S_n}((1\ 2)(3\ 4))| = 8 (n-4)!$ for all $n \ge 4$. Determine the elements in this centralizer explicitly.
\end{exercise}

\begin{sol}
	The $(n - 4)!$ factor arises from the fact that permutations on the $n - 4$ integers not involved in the cycle $(1\ 2)(3\ 4)$ commute with it. Therefore, we need only consider the permutations of $\{1, 2, 3, 4\}$ that commute with $(1\ 2)(3\ 4)$. Computing $C_{S_4}((1\ 2)(3\ 4))$ directly, the permutations that fall in this set need to leave $(1\ 2)(3\ 4)$ unchanged, or swap the two transpositions. The permutations that leave $(1\ 2)(3\ 4)$ unchanged are $1$, $(1\ 2)$, $(3\ 4)$, and $(1\ 2)(3\ 4)$. The permutations that swap the two transpositions are $(1\ 3)(2\ 4)$, $(1\ 4)(2\ 3)$, $(1\ 3\ 2\ 4)$, and $(1\ 4\ 2\ 3)$. Therefore, 
	\[C_{S_4}((1\ 2)(3\ 4)) = \{1, (1\ 2), (3\ 4), (1\ 2)(3\ 4), (1\ 3)(2\ 4), (1\ 4)(2\ 3), (1\ 3\ 2\ 4), (1\ 4\ 2\ 3)\},\]
	which has size 8. Combining this with the $(n - 4)!$ factor, we have $|C_{S_n}((1\ 2)(3\ 4))| = 8(n - 4)!$ for all $n \geq 4$.
\end{sol}

\begin{exercise}
	Let $\sigma$ be the 5-cycle $(1\ 2\ 3\ 4\ 5)$ in $S_5$. In each of (a) to (c) find an explicit element $\tau \in S_5$ which accomplishes the specified conjugation:
	\begin{subexercises}
		\item $\tau \sigma \tau^{-1} = \sigma^2$
		\item $\tau \sigma \tau^{-1} = \sigma^{-1}$
		\item $\tau \sigma \tau^{-1} = \sigma^{-2}$
	\end{subexercises}
\end{exercise}

\begin{solenum}
	\item $\sigma^2 = (1\ 3\ 5\ 2\ 4)$. Then $\tau = (2\ 3\ 5\ 4)$.
	\item $\sigma^{-1} = (1\ 5\ 4\ 3\ 2)$. Then $\tau = (1\ 5)(2\ 4)$.
	\item $\sigma^{-2} = (1\ 4\ 2\ 5\ 3)$. Then $\tau = (2\ 4\ 5\ 3)$.
\end{solenum}

\begin{exercise}
	In each of (a) -- (d) determine whether $\sigma_1$ and $\sigma_2$ are conjugate. If they are, given an explicit permutation $\tau$ such that $\tau \sigma_1 \tau^{-1} = \sigma_2$.
	\begin{subexercises}
		\item $\sigma_1 = (1\ 2)(3\ 4\ 5)$ and $\sigma_2 = (1\ 2\ 3)(4\ 5)$
		\item $\sigma_1 = (1\ 5)(3\ 7\ 2)(10\ 6\ 8\ 11)$ and $\sigma_2 = (3\ 7\ 5\ 10)(4\ 9)(13\ 11\ 2)$
		\item $\sigma_1 = (1\ 5)(3\ 7\ 2)(10\ 6\ 8\ 11)$ and $\sigma_2 = \sigma_1^3$
		\item $\sigma_1 = (1\ 3)(2\ 4\ 6)$ and $\sigma_2 = (3\ 5)(2\ 4)(5\ 6)$
	\end{subexercises}
\end{exercise}

\begin{solenum}
	\item Rewrite $\sigma_2$ as $(4\ 5)(1\ 2\ 3)$ so that both have the same cycle type. Then $\tau = (1\ 4\ 2\ 5\ 3)$.
	\item Rewrite $\sigma_2$ as $(4\ 9)(13\ 11\ 2)(3\ 7\ 5\ 10)$. Then $\tau = (1\ 4)(8\ 5\ 9)(6\ 7\ 11\ 10\ 3\ 13)$.
	\item Note that $\sigma_2 = (1\ 5)(6\ 10\ 11\ 8)$, which does not have the same cycle type as $\sigma_1$. Therefore, they are not conjugate.
	\item $\sigma_2 = (2\ 4)(3\ 5\ 6)$. Then $\tau = (1\ 2\ 3\ 4\ 5)$.
\end{solenum}

\begin{exercise}
	Find a representative for each conjugacy class of elements of order $4$ in $S_8$ and $S_{12}$.
\end{exercise}

\begin{sol}
	$S_8$:
	\[\begin{array}{c|c}
		\textbf{Cycle Type} & \textbf{Representative} \\
		\hline
		4, 4 & (1\ 2\ 3\ 4)(5\ 6\ 7\ 8) \\
		4, 2, 2 & (1\ 2\ 3\ 4)(5\ 6)(7\ 8) \\
		4, 2, 1, 1 & (1\ 2\ 3\ 4)(5\ 6) \\
		4, 1, 1, 1, 1 & (1\ 2\ 3\ 4)
	\end{array}\]
	$S_{12}$:
	\[\begin{array}{c|c}
		\textbf{Cycle Type} & \textbf{Representative} \\
		\hline
		4, 4, 4 & (1\ 2\ 3\ 4)(5\ 6\ 7\ 8)(9\ 10\ 11\ 12) \\
		4, 4, 2, 2 & (1\ 2\ 3\ 4)(5\ 6\ 7\ 8)(9\ 10)(11\ 12) \\
		4, 4, 2, 1, 1 & (1\ 2\ 3\ 4)(5\ 6\ 7\ 8)(9\ 10) \\
		4, 4, 1, 1, 1, 1 & (1\ 2\ 3\ 4)(5\ 6\ 7\ 8) \\
		4, 2, 2, 2, 2 & (1\ 2\ 3\ 4)(5\ 6)(7\ 8)(9\ 10)(11\ 12) \\
		4, 2, 2, 2, 1, 1 & (1\ 2\ 3\ 4)(5\ 6)(7\ 8)(9\ 10) \\
		4, 2, 2, 1, 1, 1, 1 & (1\ 2\ 3\ 4)(5\ 6)(7\ 8) \\
		4, 2, 1, 1, 1, 1, 1, 1 & (1\ 2\ 3\ 4)(5\ 6) \\
		4, 1, 1, 1, 1, 1, 1, 1, 1 & (1\ 2\ 3\ 4)
	\end{array} \qh\]
\end{sol}

\begin{specialexercise}
	Find all finite groups which have exactly two conjugacy classes.
\end{specialexercise}

\begin{sol}
	Let $G$ be a finite group with exactly two conjugacy classes. One of these classes must be $\{1\}$, the identity element. Let $\kk$ be the other conjugacy class, and let $g \in \kk$. By Proposition 4.6, we have $|\kk| = |G : C_G(g)|$. Since there are only two conjugacy classes, then $|\kk| = |G| - 1$. Therefore, $|G : C_G(g)| = |G| - 1$, which implies that $|C_G(g)| = |G|/(|G| - 1)$. Since $|C_G(g)|$ is an integer, then $|G| - 1$ divides $|G|$. This is only possible if $|G| - 1 = 1$, so that $|G| = 2$. Therefore, the only finite group with exactly two conjugacy classes is the group of order 2, which is isomorphic to $Z_2$.
\end{sol}

\begin{exercise}
	In \exref{4.2.1} two labellings of the elements $\set{1, a, b, c}$ of the Klein 4-group $V_4$ were chosen to give two versions of the left regular representation of $V_4$ into $S_4$. Let $\pi_1$ be the version of regular representation obtained in part (a) of that exercise, and let $\pi_2$ be the version obtained via the labelling in part (b). Let $\tau = (2\ 4)$. Show that $\tau \circ \pi_1(g) \circ \tau\inv = \pi_2(g)$ for each $g \in V_4$ (i.e., conjugation by $\tau$ sends the image of $\pi_1$ to the image of $\pi_2$ elementwise).
\end{exercise}

\begin{sol}
	We will compute for $a$, as the rest follow similarly. We have $\pi_1(a) = (1\ 2)(3\ 4)$ and $\pi_2(a) = (1\ 4)(2\ 3)$. Then $\tau \circ \pi_1(a) \circ \tau\inv = (2\ 4)(1\ 2)(3\ 4)(2\ 4) = (1\ 4)(2\ 3) = \pi_2(a)$.
\end{sol}

\begin{exercise}
	Find an element of $S_8$ which conjugates the subgroup of $S_8$ obtained in part (a) of \exref{4.2.3} to the subgroup of $S_8$ obtained in part (b) of that same exercise (both of these subgroups are isomorphic to $D_8$).
\end{exercise}

\begin{sol}
	Recall by \exref{1.7.17} that left conjugation is an automorphism, so it suffices to find an element which sends the generators of the first subgroup to the generators of the second subgroup. The generators of the first subgroup are $(1\ 2\ 3\ 4)(5\ 8\ 7\ 6)$ and $(1\ 5)(2\ 6)(3\ 7)(4\ 8)$, while the generators of the second subgroup are $(1\ 3\ 5\ 7)(2\ 8\ 6\ 4)$ and $(1\ 2)(3\ 4)(5\ 6)(7\ 8)$ respectively. It suffices to check some element $\tau$ which sends the first generator to the second; the second generator will follow automatically. One such element is $\tau = (5\ 2\ 3)(6\ 4\ 7)$.
\end{sol}

\begin{exercise}
	Find an element of $S_4$ which conjugates the subgroup of $S_4$ obtained in part (a) of \exref{4.2.5} to the subgroup of $S_4$ obtained in part (b) of that same exercise (both of these subgroups are isomorphic to $D_8$).
\end{exercise}

\begin{sol}
	We proceed similarly to the preceding exercise. The generators of the first subgroup are $(1\ 2\ 3\ 4)$ and $(2\ 4)$, while the generators of the second subgrgoup are $(1\ 3\ 2\ 4)$ and $(3\ 4)$. We obtain $\tau = (2\ 3)$.
\end{sol}

\begin{specialexercise}[4.3.17]
	Let $A$ be a nonempty set and let $X$ be any subset of $S_A$. Let
	\[F(X) = \set{a \in A \mid \sigma(a) = a \text{ for all } \sigma \in X} \qquad \text{---the \emph{fixed set} of $X$}\]
	Let $M(X) = A - F(X)$ be the elements which are \emph{moved} by some element of $X$. Let $D = \set{\sigma \in S_A \mid |M(\sigma)| < \infty}$. Prove that $D$ is a normal subgroup of $S_A$.
\end{specialexercise}

\begin{sol}
	Note that $1 \in D$, since $M(1) = \varnothing$ as it fixes all elements of $A$. Let $\sigma, \tau \in D$. Then $|M(\sigma)| < \infty$ and $|M(\tau)| < \infty$. Suppose $a \in F(\sigma) \cap F(\tau)$. Clearly, $a \in F(\sigma\tau)$. By contrapositive, then $a \in M(\sigma\tau)$ implies $a \in M(\sigma) \cup M(\tau)$, which are both finite, hence $M(\sigma\tau)$ is finite. Therefore, $\sigma\tau \in D$.
	
	Now consider $\sigma \in D$. If $a \in A$ is fixed by $\sigma$, then it is fixed by $\sigma\inv$. Therefore, any element moved by $\sigma\inv$ must be moved by $\sigma$, and vice-versa. Hence, $M(\sigma\inv) = M(\sigma)$, which is finite. Therefore, $\sigma\inv \in D$.

	Finally, consider $\sigma \in D$ and $\tau \in S_A$. Suppose $a \in F(\tau\sigma\tau\inv)$ so that $\tau\sigma\tau\inv(a) = a$. Then $\sigma(\tau\inv(a)) = \tau\inv(a)$, hence $\tau\inv(a) \in F(\sigma)$. By contrapositive, if $a \in M(\tau\sigma\tau\inv)$, then $\tau\inv(a) \in M(\sigma)$. Since $M(\sigma)$ is finite, then $M(\tau\sigma\tau\inv)$ is finite. Therefore, $\tau\sigma\tau\inv \in D$.
\end{sol}

\begin{exercise}
	Let $A$ be a set, let $H$ be a subgroup of $S_A$, and let $F(H)$ be the fixed points of $H$ on $A$ as defined in the preceding exercise. Prove that if $\tau \in N_{S_A}(H)$, then $\tau$ stabilizes the set $F(H)$ and its complement $A - F(H)$.
\end{exercise}

\begin{sol}
	Suppose $\sigma \in H$. Since $\tau \in N_{S_A}(H)$, we have $\tau\inv \in N_{S_A}(H)$, hence $\tau\inv\sigma\tau \in H$. Suppose $a \in F(H)$. Then $\tau\inv\sigma\tau(a) = a$, so that $\sigma\tau(a) = \tau(a)$. Therefore, $\tau(a) \in F(H)$, and $\tau(F(H)) \subseteq F(H)$. Bijectivity of $\tau$ also implies that $\tau\inv(F(H)) \subseteq F(H)$, which shows that $F(H) \subseteq \tau(F(H))$, hence $\tau(F(H)) = F(H)$. Therefore, $\tau$ stabilizes $F(H)$. Moreover, since $\tau$ is a bijection, it also stabilizes the complement $A - F(H)$.
\end{sol}

\begin{exercise}
	Assume $H$ is a normal subgroup of $G$, $\kk$ is a conjugacy class of $G$ contained in $H$ and $x \in \kk$. Prove that $\kk$ is a union of $k$ conjugacy classes of equal size in $H$, where $k = |G : HC_G(x)|$. Deduce that a conjugacy class in $S_n$ which consists of even permutations is either a single conjugacy class under the action of $A_n$ or is a union of two classes of the same size in $A_n$. [Let $A = C_G(x)$ and $B = H$ so $A \cap B = C_H(x)$. Draw the lattice diagram associated to the Second Isomorphism Theorem and interpret the appropriate indices. See also \exref{4.1.9}.]
\end{exercise}

\begin{sol}
	Let $G$ act on $\kk$ by conjugation. Since $\kk$ is a conjugacy class of $G$, this action is transitive. Because $H \nsub G$ and $\kk \subseteq H$, then $H$ also acts on $\kk$ by conjugation. Let $\oo$ be one such orbit of this action, and suppose $y \in \oo$. By Proposition 4.6, we have $|\oo| = |H : C_H(y)|$. By normality of $H$ in $G$, we have $C_H(y) = C_G(y) \cap H$. Using the Second Isomorphism Theorem, we have
	\[HC_G(y)/C_G(y) \cong H/C_H(y)\]
	so that $|H : C_H(y)| = |HC_G(y) : C_G(y)|$. Therefore, $|\oo| = |HC_G(y) : C_G(y)|$. Using \exref{4.1.9}, we know that each orbit of $H$ on $\kk$ has the same size, so every orbit has size $|HC_G(y) : C_G(y)|$. Suppose there are $k$ orbits. Then
	\[|\kk| = k|HC_G(y) : C_G(y)|.\]
	By Proposition 4.6, we also have $|\kk| = |G : C_G(y)|$. Therefore,
	\[|G : C_G(y)| = k|HC_G(y) : C_G(y)|.\]
	Dividing both sides by $|HC_G(y) : C_G(y)|$, we obtain $k = |G : HC_G(y)|$. Hence, $\kk$ is a union of $k$ conjugacy classes of equal size in $H$.

	Now suppose $\kk \subseteq A_n$ is a conjugacy class of $S_n$ consisting of even permutations. Let $x \in \kk$. Applying the preceding result with $H = A_n$, we have that $\kk$ is a union of $k$ conjugacy classes of equal size in $A_n$, where $k = |S_n : A_nC_{S_n}(x)|$. Since $|S_n : A_n| = 2$, then $k$ is either 1 or 2. Therefore, $\kk$ is either a single conjugacy class under the action of $A_n$ or is a union of two classes of the same size in $A_n$.
\end{sol}

\begin{exercise}
	Let $\sigma \in A_n$. Show that all elements in the conjugacy class of $\sigma$ in $S_n$ (i.e., all elements of the same cycle type as $\sigma$) are conjugate in $A_n$ if and only if $\sigma$ commutes with an odd permutation. [Use the preceding exercise.]
\end{exercise}

\begin{sol}
	\rightimp Recall by the previous exercise that the conjugacy class of $\sigma$ in $S_n$ is either a single conjugacy class in $A_n$ or a union of two conjugacy classes of equal size in $A_n$. Suppose all elements in the conjugacy class of $\sigma$ in $S_n$ are conjugate in $A_n$. Then the conjugacy class of $\sigma$ in $S_n$ is a single conjugacy class in $A_n$, so that $k = |S_n : A_nC_{S_n}(\sigma)| = 1$. Therefore, $S_n = A_nC_{S_n}(\sigma)$, so that there exists some odd permutation $\tau \in S_n$ such that $\tau \in C_{S_n}(\sigma)$. Hence, $\sigma$ commutes with an odd permutation.

	\leftimp Suppose $\sigma$ commutes with an odd permutation $\tau \in S_n$. Then $\tau \in C_{S_n}(\sigma)$, so that $A_nC_{S_n}(\sigma)$ contains odd permutations. Therefore, $|S_n : A_nC_{S_n}(\sigma)| = 1$, so that the conjugacy class of $\sigma$ in $S_n$ is a single conjugacy class in $A_n$. Hence, all elements in the conjugacy class of $\sigma$ in $S_n$ are conjugate in $A_n$.
\end{sol}

\newpage

\begin{specialexercise}
	Let $\kk$ be a conjugacy class in $S_n$ and assume that $\kk \subseteq A_n$. Show $\sigma \in S_n$ does \emph{not} commute with any odd permutation if and only if the cycle type of $\sigma$ consists of distinct odd integers. Deduce that $\kk$ consists of two conjugacy clases in $A_n$ if and only if the cycle type of an element of $\kk$ consists of distinct odd integers. [Assume first that $\sigma \in \kk$ does not commute with any odd permutation. Observe that $\sigma$ commutes with each individual cycle in its cycle decomposition---use this to show that all its cycles must be of odd length. If two cycles have the same odd length, $k$, find a product of $k$ transpositions which interchanges them and commutes with $\sigma$. Conversely, if the cycle type of $\sigma$ consists of distinct integers, prove that $\sigma$ commutes \emph{only} with the group generated by the cycles in its cycle decomposition.]
\end{specialexercise}

\begin{sol}
	\rightimp Suppose $\sigma \in S_n$ does not commute with any odd permutation. Let the cycle decomposition of $\sigma$ be $\sigma = \tau_1\tau_2\cdots\tau_k$, where each $\tau_i$ is a disjoint cycle. Since $\sigma$ commutes with each $\tau_i$, then each $\tau_i$ must be of odd length; otherwise, $\tau_i$ would be an odd permutation commuting with $\sigma$. Now suppose two cycles $\tau_i$ and $\tau_j$ have the same odd length $m$. Then the product of $m$ transpositions which interchanges the elements of $\tau_i$ and $\tau_j$ is an odd permutation which commutes with $\sigma$, contradicting our assumption. Therefore, all cycles in the cycle decomposition of $\sigma$ must have distinct odd lengths.

	\leftimp Suppose the cycle type of $\sigma$ consists of distinct odd integers. Let the cycle decomposition of $\sigma$ be $\sigma = \tau_1\tau_2\cdots\tau_k$, where each $\tau_i$ is a disjoint cycle of odd length. Any permutation that commutes with $\sigma$ must permute the cycles $\tau_i$ among themselves. However, since the lengths of the cycles are distinct, the only way to permute them while preserving their lengths is to leave them unchanged. Therefore, any permutation that commutes with $\sigma$ must be a product of powers of the individual cycles $\tau_i$. Since each $\tau_i$ has odd length, any such product will be an even permutation. Hence, $\sigma$ does not commute with any odd permutation.
\end{sol}

\begin{exercise}
	Show that if $n$ is odd, then the set of all $n$-cycles consists of two conjugacy classes of equal size in $A_n$.
\end{exercise}

\begin{sol}
	Let $\sigma \in S_n$ be an $n$-cycle. Since $n$ is odd, the cycle type of $\sigma$ consists of the single odd integer $n$, which is distinct. By the previous exercise, $\sigma$ does not commute with any odd permutation. Therefore, by the exercise before that, the conjugacy class of $\sigma$ in $S_n$ consists of two conjugacy classes of equal size in $A_n$. Since this holds for any $n$-cycle $\sigma$, the set of all $n$-cycles consists of two conjugacy classes of equal size in $A_n$.
\end{sol}

\begin{exercise}[4.3.23]
	Recall (cf. \exref{2.4.16}) that a proper subgroup $M$ of $G$ is called \emph{maximal} if whenever $M \leq H \leq G$, either $H = M$ or $H = G$. Prove that if $M$ is a maximal subgroup of $G$ then either $N_G(M) = M$ or $N_G(M) = G$. Deduce that if $M$ is a maximal subgroup of $G$ that is not normal in $G$, then the number of nonidentity elements of $G$ that are contained in conjugates of $M$ is at most $(|M| - 1)|G : M|$.
\end{exercise}

\begin{sol}
	Suppose $M$ is a maximal subgroup of $G$. By definition, $N_G(M)$ is a subgroup of $G$ containing $M$. Therefore, by maximality of $M$, either $N_G(M) = M$ or $N_G(M) = G$.

	Now suppose $M$ is a maximal subgroup of $G$ that is not normal in $G$. Then by the previous result, we have $N_G(M) = M$. By Proposition 4.6, the size of the conjugacy class of $M$ in $G$ is $|G : N_G(M)| = |G : M|$. Each conjugate of $M$ contains $|M| - 1$ nonidentity elements. Therefore, the total number of nonidentity elements of $G$ contained in conjugates of $M$ is at most $(|M| - 1)|G : M|$.
\end{sol}

\newpage

\begin{exercise}[4.3.24]
	Assume $H$ is a proper subgroup of the finite group $G$. Prove
	\[G \neq \bigcup_{g \in G} gHg\inv\]
	i.e., $G$ is not the union of the conjugates of any proper subgroup. [Put $H$ in some maximal subgroup and use the preceding exercise.]
\end{exercise}

\begin{sol}
	Suppose $H$ is a proper subgroup of the finite group $G$. Then there exists a maximal subgroup $M$ of $G$ such that $H \leq M < G$. By the previous exercise, either $N_G(M) = M$ or $N_G(M) = G$, so that we have two cases: either $M \nsub G$ or $M$ is not normal in $G$.

	If $M$ is normal in $G$, then all conjugates of $M$ are equal to $M$. Therefore, the union of the conjugates of $H$ is contained in $M$, which is a proper subset of $G$. Hence, 
	\[G \neq \bigcup_{g \in G} gHg\inv.\]
	
	
	Now suppose $M$ is not normal in $G$. Then by the previous exercise, the number of nonidentity elements of $G$ contained in conjugates of $M$ is at most $(|M| - 1)|G : M|$. Moreover, $H \leq M$ implies that any conjugate of $H$ is contained in a conjugate of $M$. Therefore, the number of nonidentity elements of $G$ contained in conjugates of $H$ is also at most $(|M| - 1)|G : M|$. Since $M$ is a proper subgroup of $G$, then $|G : M| \geq 2$, so that $(|M| - 1)|G : M| < |G| - 1$. Therefore, there exists at least one nonidentity element of $G$ that is not contained in any conjugate of $H$, hence
	\[G \neq \bigcup_{g \in G} gHg\inv. \qh\]
\end{sol}

\begin{exercise}[4.3.25]
	Let $G = \gl_2(\c)$ and let
	\[H = \set*{\left.
		\begin{pmatrix}
			a & b \\
			0 & c
		\end{pmatrix}~\right| a, b, c \in \c, ac \neq 0
	}.\]
	Prove that every element of $G$ is conjugate to some element of the subgroup $H$ and deduce that $G$ is the union of conjugates of $H$. [Show that every element of $\gl_2(\c)$ has an eigenvector.]
\end{exercise}

\begin{sol}
	Suppose $X \in G$, where $X$ is the matrix
	\[
	\begin{pmatrix}
		a & b \\
		c & d
	\end{pmatrix}\]
	and consider the characteristic polynomial $\chi_X(t) = \det(X - tI) = t^2 - (a + d)t + (ad - bc)$. Since $\c$ is closed under complex multiplication, $\chi_X(t)$ has at least one root $\lambda \in \c$. Therefore, there exists a nonzero vector $\b v \in \c^2$ such that $X\b v = \lambda \b v$, so that $\b v$ is an eigenvector of $X$ corresponding to the eigenvalue $\lambda$. We may extend $\b v$ to a basis $\set{\b v, \b w}$ of $\c^2$. Let $P$ be the matrix whose columns are $\b v$ and $\b w$. Then $P$ is invertible, and we have
	\[P\inv XP = \begin{pmatrix}
		\lambda & \alpha \\
		0 & \beta
	\end{pmatrix} \in H.\]
	for some $\alpha, \beta \in \c$ such that $X\b w = \alpha \b v + \beta \b w$. Therefore, every element of $G$ is conjugate to some element of the subgroup $H$, hence every element of $G$ is contained in some conjugate of $H$. Thus,
	\[G = \bigcup_{g \in G} gHg\inv. \qh\]
\end{sol}

\begin{exercise}
	Let $G$ be a transitive permutation group on the finite set $A$ with $\abs A > 1$. Show that there is some $\sigma \in G$ such that $\sigma(a) \neq a$ for all $a \in A$ (such an element $\sigma$ is called \emph{fixed point free}).
\end{exercise}

\begin{sol}
	For each $a \in A$, consider $G_a$. By transitivity of $G$ on $A$, we have $|G : G_a| = |A| > 1$, so that $G_a$ is a proper subgroup of $G$. By \exref{4.3.24}, we have
	\[G \neq \bigcup_{g \in G} gG_ag\inv.\]
	Therefore, there exists some $\sigma \in G$ such that $\sigma \notin gG_ag\inv$ for all $g \in G$. Moreover, if $\sigma(a) = a$ for some $a \in A$, then $\sigma \in G_a$, hence $\sigma \in gG_ag\inv$ for $g = 1$, a contradiction. Therefore, $\sigma(a) \neq a$ for all $a \in A$.
\end{sol}

\begin{exercise}
	Let $g_1, g_2, \ldots, g_r$ be representatives of the conjugacy classes of the finite group $G$ and assume these elements pairwise commute. Prove that $G$ is abelian.
\end{exercise}

\begin{sol}
	It is clear that $g_i \in C_G(g_j)$ for any two representatives $g_i$ and $g_j$. Let $S$ be the set of class representatives, so that we may write $S \subseteq C_G(g_i)$. Then $|C_G(g_i)| \geq |S| = r$ for all $i$. From the class equation, we get that
	\[|G| = \sum_{i = 1}^r |G : C_G(g_i)| \leq \sum_{i = 1}^r \frac{|G|}{r} = |G|\]
	Each term in the sum must be equal to $|G|/r$ so that $|C_G(g_i)| = r$ for all $i$. Therefore, $C_G(g_i) = S$ for all $i$. Recall that one of the $g_i$ must be 1, so that $C_G(1) = G$. Therefore, $G = S$, so that $G$ is abelian.
\end{sol}

\begin{exercise}
	Let $p$ and $q$ be primes with $p < q$. Prove that a non-abelian group $G$ of order $pq$ has a nonnormal subgroup of index $q$ so that there exists an injective homomorphism into $S_q$. Deduce that $G$ is isomorphic to a subgroup of the normalizer in $S_q$ of the cyclic group generated by the $q$-cycle $(1\ 2\ \cdots\ q)$.
\end{exercise}

\begin{sol}
	By Cauchy's Theorem, there exists $g \in G$ such that $\abs g = q$, hence $\gen g$ has order $q$ and is of index $p$. Moreover, $\gen g$ is not normal in $G$, for otherwise $G/\gen g$ would be cyclic and isomorphic to $Z_p$, contradicting that $G$ is non-abelian. Similarly, $G$ also contains a subgroup $P$ of order $p$ and index $q$ and is similarly not normal in $G$.

	Let $G$ act on the left cosets of $P$ in $G$ by left multiplication. Then there is the associated permutation representation $\pi : G \to S_q$. In particular, $\pi$ is injective: If $\pi(x) = 1$ for some $x \in G$, then $xgP = gP$ for every $g \in G$, or $g\inv xg \in P$ for every $g \in G$. But $\ker\pi \nsub G$ and is contained in $P$ so that $\ker\pi$ has order 1 or $p$. If $|\ker\pi| = p$, then $\ker\pi = P$ is normal in $G$, a contradiction. Therefore, $\ker\pi = 1$ and $\pi$ is injective.

	Recall that $\abs g = q$ so that $\abs{\pi(g)} = q$, so $\pi(g)$ acts transitively on $G/P$ since its powers generate $q$ distinct cosets. Now pick $h \in G$. From \exref{1.1.22}, it follows that $|hgh\inv| = |g| = q$ so that $hgh\inv = g^k$ for some integer $k$ with $1 \leq k < q$. Therefore,
	\[\pi(h)\pi(g)\pi(h)\inv = \pi(hgh\inv) = \pi(g^k) = (\pi(g))^k \in \gen{\pi(g)}.\]
	Hence, $\pi(h) \in N_{S_q}(\gen{\pi(g)})$. Since $h$ was arbitrary, we have $G \cong \pi(G) \leq N_{S_q}(\gen{\pi(g)})$. Therefore, $G$ is isomorphic to a subgroup of the normalizer in $S_q$ of the cyclic group generated by the $q$-cycle $(1\ 2\ \cdots\ q)$.
\end{sol}

\begin{exercise}
	Let $p$ be a prime and let $G$ be a group of order $p^\alpha$. Prove that $G$ has a subgroup of order $p^\beta$, for every $\beta$ with $0 \leq \beta \leq \alpha$. [Use Theorem 8 and induction on $\alpha$.]
\end{exercise}

\begin{sol}
	We proceed by induction on $\alpha$. If $\alpha = 0$, then $G$ is the trivial group, which has a subgroup of order $p^0 = 1$. Now suppose the result holds for all groups of order $p^k$ for some $k \geq 0$. Let $G$ be a group of order $p^{k + 1}$. By Theorem 8, $Z(G)$ is nontrivial, so there exists some $x \in Z(G)$ such that $\abs x = p$. Then $\gen x$ is a normal subgroup of $G$ of order $p$. Consider the quotient group $G/\gen x$, which has order $p^k$. By the induction hypothesis, for every $\beta$ with $0 \leq \beta \leq k$, there exists a subgroup $H/\gen x$ of $G/\gen x$ such that $\abs{H/\gen x} = p^\beta$. By the Lattice Isomorphism Theorem, there exists a subgroup $H$ of $G$ such that $\gen x \leq H$ and $\abs{H} = p^{\beta + 1}$. Therefore, for every $\beta$ with $1 \leq \beta \leq k + 1$, there exists a subgroup of $G$ of order $p^\beta$. Since the trivial subgroup has order $p^0 = 1$, the result holds for all $\beta$ with $0 \leq \beta \leq k + 1$. By induction, the result holds for all $\alpha \geq 0$.
\end{sol}

\begin{exercise}
	If $G$ is a group of odd order, prove for any nonidentity element $x \in G$ that $x$ and $x^{-1}$ are not conjugate in $G$.
\end{exercise}

\begin{sol}
	Suppose for contradiction that $x$ and $x\inv$ are conjugate in $G$. Then there exists $g \in G$ such that $gxg\inv = x\inv$. If we conjugate both sides by $g$ again, we obtain
	\[g^2xg^{-2} = gx\inv g\inv = (gxg\inv)\inv = (x\inv)\inv = x\]
	so that $g^2x = xg^2$. Then $g^2 \in C_G(x)$. Consider the quotient group $G/C_G(x)$. Since $g^2 \in C_G(x)$, we have $(gC_G(x))^2 = C_G(x)$, so that the order of $gC_G(x)$ in $G/C_G(x)$ is 1 or 2. However, since $G$ has odd order, then $G/C_G(x)$ also has odd order, so the order of $gC_G(x)$ cannot be 2. Therefore, the order of $gC_G(x)$ is 1, so that $g \in C_G(x)$. But this implies that $gxg\inv = x$, contradicting our assumption. Therefore, $x$ and $x\inv$ are not conjugate in $G$.
\end{sol}

\begin{exercise}
	Using the usual generators and relations for the dihedral group $D_{2n}$ (cf. Section 1.2) show that for $n = 2k$ an even integer the conjugacy classes in $D_{2n}$ are the following: $\set 1, \set{r^k}, \set{r^{\pm 1}}, \set{r^{\pm 2}}, \ldots, \set{r^{\pm(k - 1)}}, \set{sr^{2b} \mid b = 1, \ldots, k}$ and $\set{sr^{2b - 1} \mid b = 1, \ldots, k}$. Give the class equation for $D_{2n}$.
\end{exercise}

\begin{sol}
	We immediately have two conjugacy classes, $\set 1$ and $\set{r^n}$, since $r^n \in Z(D_{2n})$ when $n$ is even. Moreover, recall that the elements of $D_{2n}$ are of the form $r^m$ or $sr^m$ for some integer $m$. Consider the conjugacy class of $r^k$ for some $1 \leq k < n$ except $k = n/2$. For any $j$, we have
	\[r^jr^kr^{-j} = r^k, \quad \text{and} \quad sr^jr^kr^{-j}s = r^{-k}\]
	so that the conjugacy class of $r^k$ is $\set{r^{\pm k}}$. Now consider the conjugacy class of $sr^m$ for some integer $m$. For any $j$, we have
	\[r^jsr^mr^{-j} = sr^{m - 2j}, \quad \text{and} \quad sr^jsr^mr^{-j}s = r^{-(m - 2j)}\]
	so that the conjugacy class of $sr^m$ is $\set{sr^{m - 2j} \mid j \in \z}$. If $m$ is even, then this conjugacy class is $\set{sr^{2b} \mid b = 1, \ldots, k}$; if $m$ is odd, then this conjugacy class is $\set{sr^{2b - 1} \mid b = 1, \ldots, k}$. Therefore, the conjugacy classes in $D_{2n}$ when $n$ is even are $\set 1, \set{r^k}, \set{r^{\pm 1}}, \set{r^{\pm 2}}, \ldots, \set{r^{\pm(k - 1)}}, \set{sr^{2b} \mid b = 1, \ldots, k}$ and $\set{sr^{2b - 1} \mid b = 1, \ldots, k}$. Lastly, the class equation for $D_{2n}$ is
	\[|D_{2n}| = 1 + 1 + 2(k - 1) + k + k.\qedhere\]
\end{sol}

\begin{exercise}
	For $n = 2k + 1$ an odd integer, show that the conjugacy classes in $D_{2n}$ are $\set{1}, \set{r^{\pm 1}}, \set{r^{\pm 2}}, \ldots, \set{r^{\pm k}}$, and $\set{sr^b \mid b = 1, \ldots, n}$. Give the class equation for $D_{2n}$.
\end{exercise}

\begin{sol}
	We immediately have the conjugacy class $\set 1$. Moreover, recall that the elements of $D_{2n}$ are of the form $r^m$ or $sr^m$ for some integer $m$. Consider the conjugacy class of $r^k$ for some $1 \leq k \leq n$ except $k = n/2$. For any $j$, we have
	\[r^jr^kr^{-j} = r^k, \quad \text{and} \quad sr^jr^kr^{-j}s = r^{-k}\]
	so that the conjugacy class of $r^k$ is $\set{r^{\pm k}}$. Now consider the conjugacy class of $sr^m$ for some integer $m$. For any $j$, we have
	\[r^jsr^mr^{-j} = sr^{m - 2j}, \quad \text{and} \quad sr^jsr^mr^{-j}s = r^{-(m - 2j)}\]
	so that the conjugacy class of $sr^m$ is $\set{sr^{m - 2j} \mid j \in \z}$. Since $n$ is odd, this conjugacy class is $\set{sr^b \mid b = 1, \ldots, n}$. Therefore, the conjugacy classes in $D_{2n}$ when $n$ is odd are $\set{1}, \set{r^{\pm 1}}, \set{r^{\pm 2}}, \ldots, \set{r^{\pm k}}$, and $\set{sr^b \mid b = 1, \ldots, n}$. Lastly, the class equation for $D_{2n}$ is
	\[|D_{2n}| = 1 + 2k + n.\qedhere\]
\end{sol}

\begin{specialexercise} [4.3.33]
	This exercise gives a formula for the size of each conjugacy class in $S_n$. Let $\sigma$ be a permutation in $S_n$ and let $m_1, m_2, \ldots, m_s$ be the \emph{distinct} integers which appear in the cycle type of $\sigma$ (including 1-cycles). For each $i \in \set{1, 2, \ldots, s}$ assume $\sigma$ has $k_i$ cycles of length $m_i$ (so that $\sum_1^s k_im_i = n$). Prove that the number of conjugates of $\sigma$ is
	\[\frac{n!}{(k_1!m_1^{k_1})(k_2!m_2^{k_2}) \cdots (k_s!m_s^{k_s})}\]
	[See \exref{1.3.6} and \exref{1.3.7} where this formula was given in some special cases.]
\end{specialexercise}

\begin{sol}
	We start with $n$ integers. We see that there are $n!$ ways to arrange these integers. Now consider the $k_1$ cycles of length $m_1$ in the cycle decomposition of $\sigma$. Clearly, we must use $k_1m_1$ of the $n$ integers to form these cycles. However, there are $m_1^{k_1}$ ways to arrange the integers within each of these $k_1$ cycles, and there are $k_1!$ ways to arrange the $k_1$ cycles among themselves. Therefore, we must divide $n!$ by $(k_1!m_1^{k_1})$ to account for these arrangements. Continuing in this manner for each $i$ from 1 to $s$, we see that the total number of distinct arrangements of the $n$ integers that correspond to the same cycle type as $\sigma$ is given by
	\[\frac{n!}{(k_1!m_1^{k_1})(k_2!m_2^{k_2}) \cdots (k_s!m_s^{k_s})}.\]
	Since the conjugacy class of $\sigma$ in $S_n$ consists of all permutations with the same cycle type as $\sigma$, the number of conjugates of $\sigma$ is given by this formula.
\end{sol}

\begin{exercise}
	[4.3.34] Prove that if $p$ is a prime and $P$ is a subgroup of $S_p$ of order $p$, then $|N_{S_p}(P)| = p(p - 1)$. [Argue that every conjugate of $P$ contains exactly $p - 1$ $p$-cycles and use the formula for the number of $p$-cycles to compute the index of $N_{S_p}(P)$ in $S_p$.]
\end{exercise}

\begin{sol}
	Note that $P$ is cyclic of order $p$, so $P = \gen\sigma$ for some $p$-cycle $\sigma \in S_p$. Moreover, every nonidentity element of $P$ is a $p$-cycle, so $P$ contains exactly $p - 1$ $p$-cycles. Now for some $\tau \in S_p$, the conjugate $\tau P\tau\inv = \gen{\tau\sigma\tau\inv}$ also has order $p$ and contains exactly $p - 1$ $p$-cycles. Therefore, each conjugate of $P$ contains exactly $p - 1$ $p$-cycles. Note that the number of distinct $p$-cycles in $S_p$ is given by $p!/p = (p - 1)!$.

	Recall by Proposition 4.6 that the size of the conjugacy class of $P$ in $S_p$ is given by $|S_p : N_{S_p}(P)|$. Let this index be $k$. Since each conjugate of $P$ contains exactly $p - 1$ $p$-cycles, the total number of distinct $p$-cycles contained in all conjugates of $P$ is $k(p - 1)$. However, since every $p$-cycle in $S_p$ is contained in some conjugate of $P$, we have $k(p - 1) = (p - 1)!$. Therefore, $k = (p - 2)!$, so that
	\[|N_{S_p}(P)| = \frac{|S_p|}{k} = \frac{p!}{(p - 2)!} = p(p - 1).\qedhere\]
\end{sol}

\begin{exercise}
	Let $p$ be a prime. Find a formula for the number of conjugacy classes of elements of order $p$ in $S_n$ (using the greatest integer function).
\end{exercise}

\begin{sol}
	Note that elements of order $p$ in $S_n$ are products of disjoint $p$-cycles. Moreover, each $p$-cycle is conjugate to any other $p$-cycle in $S_n$. Therefore, the conjugacy class of an element of order $p$ in $S_n$ is determined by the number of disjoint $p$-cycles in its cycle decomposition. Let $k$ be the number of disjoint $p$-cycles in the cycle decomposition of such an element. Then we have $1 \leq k \leq \floor{n/p}$, hence the number of conjugacy classes of elements of order $p$ in $S_n$ is given by $\floor{n/p}$.
\end{sol}

\begin{specialexercise}
	Let $\pi : G \to S_G$ be the left regular representation afforded by the action of $G$ on itself by left multiplication. For each $g \in G$ denote the permutation $\pi(g)$ by $\sigma_g$ so that $\sigma_g(x) = gx$ for all $x \in G$. Let $\lambda : G \to S_G$ be the permutation representation afforded by the corresponding right action of $G$ on itself, and for each $h \in G$ denote the permutation $\lambda(h)$ by $\tau_h$. Thus $\tau_h(x) = xh\inv$ for all $x \in G$ ($\lambda$ is called the \emph{right regular representation} of $G$).
	\begin{subexercises}
		\item Prove that $\sigma_g$ and $\tau_h$ commute for all $g, h \in G$. (Thus the centralizer in $S_G$ of $\pi(G)$ contains the subgroup $\lambda(G)$ which is isomorphic to $G$).
		\item Prove that $\sigma_g = \tau_g$ if and only if $g$ is an element of order 1 or 2 in the center of $G$.
		\item Prove that $\sigma_g = \tau_h$ if and only if $g$ and $h$ lie in the center of $G$. Deduce that $\pi(G) \cap \lambda(G) = \pi(Z(G)) = \lambda(Z(G))$.
	\end{subexercises}
\end{specialexercise}

\begin{solenum}
	\item For any $x \in G$, we have
	\[\sigma_g(\tau_h(x)) = \sigma_g(xh\inv) = g(xh\inv) = (gx)h\inv = \tau_h(gx) = \tau_h(\sigma_g(x)).\]
	Therefore, $\sigma_g$ and $\tau_h$ commute for all $g, h \in G$.
	\item If $\sigma_g = \tau_g$, then $gx = xg\inv$ for all $x \in G$. In particular, $x = 1$ implies $g = g\inv$ so that $g^2 = 1$, hence $\abs g$ is 1 or 2. Moreover, for any $x \in G$, we have $gx = xg\inv = xg$, so that $g \in Z(G)$
	
	If $g$ is an element of order 1 or 2 in the center of $G$, then for any $x \in G$, we have $\sigma_g(x) = gx = xg = xg\inv = \tau_g(x)$. Therefore, $\sigma_g = \tau_g$.
	\item If $\sigma_g = \tau_h$, then $gx = xh\inv$ for all $x \in G$. In particular, $x = 1$ implies $g = h\inv$, so that $h = g\inv$. Therefore, $gx = xg\inv$ for all $x \in G$, so that by part (b), $g$ lies in the center of $G$. Since $h = g\inv$, then $h$ also lies in the center of $G$.
	
	If $g$ and $h$ lie in the center of $G$, then for any $x \in G$, we have $\sigma_g(x) = gx = xg = xh\inv = \tau_h(x)$. Therefore, $\sigma_g = \tau_h$. Moreover, if $\sigma_g \in \pi(G) \cap \lambda(G)$, then there exists $h \in G$ such that $\sigma_g = \tau_h$. By the above result, both $g$ and $h$ lie in the center of $G$, so that $\sigma_g \in \pi(Z(G))$ and $\tau_h \in \lambda(Z(G))$. Conversely, if $g \in Z(G)$, then by part (b), we have $\sigma_g = \tau_g$, so that $\sigma_g \in \pi(G) \cap \lambda(G)$. Therefore, $\pi(G) \cap \lambda(G) = \pi(Z(G)) = \lambda(Z(G))$.
\end{solenum}